% Compile with XeLaTeX (Korean via kotex). Benchmark-track framing of the same study.
\documentclass[10pt,conference]{IEEEtran}

\usepackage{kotex}
\usepackage{cite}
\usepackage{amsmath,amssymb,amsfonts}
\usepackage{graphicx}
\usepackage{booktabs}
\usepackage{textcomp}
\usepackage{url}
\usepackage{listings}

\lstset{basicstyle=\ttfamily\footnotesize,breaklines=true,frame=single,columns=fullflexible,keepspaces=true}

\begin{document}

\title{AIBugBench: data leakage 없이 사람과 짝지은 LLM 생성 코드 결함 벤치마크}

\maketitle

\begin{abstract}
LLM(Large Language Model) 코드를 평가하는 기존 벤치마크는 두 가지 한계를 동시에 갖는다. 첫째, HumanEval과 MBPP 같은 공개 벤치마크는 학습 데이터에 포함되어 data leakage 우려가 크다. 둘째, 사람 버그를 포함한 데이터셋은 대부분 LLM 대중화 이전에 구축되어 생성 모델 입장에서 leakage가 있거나 AI 생성 버그를 담지 않으며, AI와 사람 버그를 같은 문제에서 짝지어 비교하기 어렵다. 본 논문은 이 공백을 메우는 벤치마크 AIBugBench를 제시한다. 핵심 설계는 모델 cutoff(\texttt{gpt-3.5-turbo}, 2021년 9월)와 ChatGPT 출시(2022년 11월) 사이의 겹치는 구간을 선택하여, 대상 문제는 모델이 학습하지 못했고(leakage 없음) 그 시기의 사람 제출은 ChatGPT 이전이라 진정함을 동시에 보장하는 것이다. AIBugBench는 ConDefects를 토대로 428개 AtCoder 문제 각각에 대해 정규화된 자연어 명세, 전체 hidden test(문제당 평균 약 32개), \texttt{gpt-3.5-turbo}가 생성한 10개의 솔루션과 채점 라벨, 그리고 같은 문제에서 짝지어진 사람 정답 및 사람 버그(fault location 포함)를 제공한다. 우리는 벤치마크가 지원하는 다섯 가지 평가 태스크(bug classification, fault localization, automatic program repair, AI 대 사람 버그 비교, spec-grounded repair)와 측정 지표를 정의하고, 생성 모델의 난이도별 통과율과 결함 분포를 reference baseline으로 보고한다. 특징 분석에서 결함의 76\%가 Spec-Misinterpretation이고 한 문제에 대한 버그가 같은 유형으로 강하게 수렴함(modal-share 0.845)을 확인하였다. 데이터셋, 코드, 평가 스크립트를 공개한다.
\end{abstract}

\begin{IEEEkeywords}
benchmark, LLM-generated code, bug dataset, data leakage, automatic program repair, fault localization
\end{IEEEkeywords}

\section{Introduction}
LLM(Large Language Model)이 작성한 코드가 교육과 산업 현장에 빠르게 들어오면서, 그 코드의 결함을 체계적으로 연구하고 수리 도구를 평가할 데이터의 필요성이 커졌다~\cite{chen2021codex,tambon2024bugs}. 그러나 현존하는 벤치마크는 이 목적에 잘 맞지 않는다. HumanEval~\cite{chen2021codex}과 MBPP~\cite{austin2021mbpp}는 손으로 만든 합성 과제로 사람 버그가 없고 널리 contamination되어 있으며~\cite{golchin2024contamination,liu2023evalplus}, CodeNet~\cite{puri2021codenet}과 Codeforces 계열은 사람 제출은 풍부하나 AI 생성 버그가 없고 LLM 이전에 구축되어 생성 모델에 leakage가 있다. 결함 위치추적과 수리용으로 leakage를 통제한 ConDefects~\cite{wu2023condefects}조차 자연어 명세와 AI 생성 버그를 포함하지 않는다.

벤치마크를 깨끗하게 만들려면 두 제약을 동시에 만족해야 한다. AI 생성이 진짜 추론의 산물이려면 대상 문제가 모델 cutoff 이후여야 하고(leakage 없음), 비교 대상 사람 코드가 진정하려면 LLM 대중화 이전에 작성되어야 한다. 전자는 ``이후''를, 후자는 ``이전''을 요구하여 충돌하는 것처럼 보인다. 본 논문의 관찰은 두 제약이 서로 다른 산출물에 서로 다른 cutoff로 걸린다는 것이다. leakage는 생성 대상 문제에, 진정성은 사람 제출에 걸리며 두 cutoff가 다르므로 그 사이에 겹치는 구간이 존재한다. AIBugBench는 \texttt{gpt-3.5-turbo} cutoff(2021년 9월)와 ChatGPT 출시(2022년 11월 30일) 사이 구간의 AtCoder 문제만 사용하여 leakage-free와 human-authentic을 동시에 달성한다.

본 논문의 기여는 다음과 같다. 첫째, leakage가 없고 같은 문제에서 사람과 짝지어진 LLM 코드 결함 벤치마크 AIBugBench를 구축한다. 둘째, 벤치마크가 지원하는 다섯 평가 태스크와 측정 지표, 표준 데이터 분할과 평가 스크립트로 이루어진 평가 프로토콜을 정의한다. 셋째, 생성 모델의 난이도별 통과율과 데이터 기반 결함 분류 분포, 그리고 한 문제 안에서 버그가 같은 유형으로 수렴하는 systematicity를 reference 특성으로 보고한다. 넷째, 데이터셋과 코드, 평가 스크립트를 재현 가능하게 공개한다.

\section{Comparison with Existing Benchmarks}
표~\ref{tab:compare}는 AIBugBench를 대표적 벤치마크와 비교한 것이다. 핵심 차별점은 AIBugBench가 (i) AI 생성 버그와 (ii) 같은 문제의 사람 버그를, (iii) leakage가 없는 동시에 사람이 진정한 구간에서, (iv) 전체 hidden test, (v) 사람 fault location, (vi) 자연어 명세와 함께 제공하는 유일한 벤치마크라는 점이다. HumanEval과 MBPP는 사람 버그가 없고 contamination이 보고되었으며, CodeNet은 사람 버그는 있으나 AI 버그가 없고 생성 모델에 leakage가 있으며, ConDefects는 leakage는 통제했으나 AI 버그와 명세가 없다.

\begin{table}[t]
\caption{Comparison of code benchmarks. AI = AI-generated bugs; Hum = paired human bugs; LF = leakage-free for the generator; HA = human-authentic; FT = full hidden tests; FL = fault locations; Spec = NL specification.}
\label{tab:compare}
\centering
\footnotesize
\begin{tabular}{lccccccc}
\toprule
Benchmark & AI & Hum & LF & HA & FT & FL & Spec \\
\midrule
HumanEval~\cite{chen2021codex} & -- & -- & -- & -- & $\sim$ & -- & \checkmark \\
MBPP~\cite{austin2021mbpp}     & -- & -- & -- & -- & $\sim$ & -- & \checkmark \\
EvalPlus~\cite{liu2023evalplus} & -- & -- & -- & -- & \checkmark & -- & \checkmark \\
CodeNet~\cite{puri2021codenet} & -- & \checkmark & -- & \checkmark & $\sim$ & -- & $\sim$ \\
ConDefects~\cite{wu2023condefects} & -- & \checkmark & \checkmark & -- & \checkmark & \checkmark & -- \\
\textbf{AIBugBench (ours)} & \checkmark & \checkmark & \checkmark & \checkmark & \checkmark & \checkmark & \checkmark \\
\bottomrule
\end{tabular}
\end{table}

\section{Benchmark Construction}
\subsection{Cutoff--Release Window}
생성 모델로 \texttt{gpt-3.5-turbo}(knowledge cutoff 2021년 9월)를 고정한다. leakage가 없으려면 문제가 2021년 9월 이후 출제여야 하고, 사람 대조군이 진정하려면 ChatGPT 출시(2022년 11월 30일) 이전 컨테스트여야 한다. 따라서 [2021-09, 2022-11] 구간만 사용한다.

\subsection{Base Data and Filtering}
ConDefects~\cite{wu2023condefects}에서 위 구간의 Python 과제 중 사람 정답과 버그 짝, 전체 hidden test를 갖춘 문제를 선별한다. AtCoder가 표시 라벨을 ``Ex''로 바꾼 마지막 문제는 ConDefects의 \texttt{\_h} 표기와 매핑하여 테스트를 복구하였고, stdin/stdout 채점이 불가능한 interactive 문제는 제외하였다. 최종적으로 428개 문제와 1{,}215개의 짝지어진 사람 버그(fault location 포함)를 얻는다.

\subsection{Specifications and Tests}
ConDefects는 자연어 명세를 포함하지 않으므로, 각 과제의 영어 명세를 AtCoder에서 수집해 구조적으로 정규화하여 \texttt{problem.txt}로 제공한다. 채점에는 ConDefects의 전체 hidden test(문제당 평균 약 32개)를 사용한다.

\subsection{AI Bug Generation}
각 문제에 대해 \texttt{gpt-3.5-turbo}로 temperature 1.0, 문제당 10개의 솔루션을 생성한다. 프롬프트는 ``Output ONLY a complete Python 3 program''이라는 system 지시와 정규화 명세(user 메시지)로 구성한다. 이 지시가 없으면 생성물의 46\%가 비-Python으로 이탈하므로, 본 설정에서 모든 생성물은 100\% Python이다. 생성물은 전체 hidden test로 채점하여 pass와 WA(Wrong Answer), TLE(Time Limit Exceeded), RE(Runtime Error), CE(Compile Error)로 라벨링한다.

\subsection{Quality Control}
공개 sample만으로 채점하면 통과율이 23\%로 과대평가되지만 전체 hidden test로는 13.8\%로 떨어지므로, 모든 채점에 전체 테스트를 사용한다. 채점 샌드박스에는 \texttt{sympy}와 \texttt{sortedcontainers} 등 표준 라이브러리를 설치하여 라이브러리 부재로 인한 거짓 RE를 방지한다. 분류 라벨은 범주 확정 후 두 저자의 독립 분류로 Cohen's $\kappa$~\cite{cohen1960kappa}를 측정하여 신뢰도를 보고한다(목표 $\kappa\geq 0.7$).

\section{Benchmark Statistics}
표~\ref{tab:stats}는 벤치마크 구성이다. 428개 문제에 대해 AI 생성물 4{,}280개(문제당 10개)와 사람 버그 1{,}215개를 제공하며, AI 생성물 중 3{,}691개가 버그이다. 모든 문제는 정규화 명세, 전체 hidden test, 사람 정답과 fault location을 갖는다.

\begin{table}[t]
\caption{Composition of AIBugBench}
\label{tab:stats}
\centering
\begin{tabular}{lr}
\toprule
Item & Count \\
\midrule
Problems (AtCoder, 2021-10--2022-11) & 428 \\
Hidden tests per problem (avg.) & $\approx$32 \\
AI generations (\texttt{gpt-3.5-turbo}, 10/problem) & 4{,}280 \\
\quad of which buggy & 3{,}691 \\
Paired human faulty programs & 1{,}215 \\
Human correct programs (reference) & provided \\
NL specifications & 428 \\
Fault locations (human) & provided \\
\bottomrule
\end{tabular}
\end{table}

표~\ref{tab:difficulty}는 생성 모델의 난이도별 통과율로, 벤치마크의 난이도 분포와 reference baseline을 함께 보여준다. 통과율은 쉬운 문제에서 45.8\%, 어려운 문제에서 사실상 0\%로, 폭넓은 난이도를 포괄한다.

\begin{table}[t]
\caption{Reference pass rate of \texttt{gpt-3.5-turbo} by difficulty (full hidden tests)}
\label{tab:difficulty}
\centering
\begin{tabular}{lrrr}
\toprule
Difficulty (AtCoder) & Gen. & Pass & Rate \\
\midrule
Easy ($<$400)        & 1{,}200 & 550 & 45.8\% \\
Medium (400--1200)   & 1{,}040 & 39  & 3.7\% \\
Hard (1200--2000)    & 890     & 0   & 0.0\% \\
Very hard ($\geq$2000) & 1{,}150 & 0 & 0.0\% \\
\midrule
Total                & 4{,}280 & 589 & 13.8\% \\
\bottomrule
\end{tabular}
\end{table}

표~\ref{tab:taxonomy}는 결함 분류 분포로 벤치마크의 결함 특성을 보여준다. Spec-Misinterpretation이 76\%로 지배적이며 Hallucinated-API는 사실상 없다. 관찰된 \texttt{ModuleNotFoundError}는 환각이 아니라 미설치 라이브러리(\texttt{sympy} 등) 때문이었다. 추가로 한 문제에 대한 버그의 systematicity를 측정한 결과, 버그가 2개 이상인 387개 문제에서 평균 modal-share가 0.845, 완전 일치 문제 비율이 38.8\%였다. 즉 LLM은 한 문제를 일관되게 같은 방식으로 오해하며, 이는 벤치마크가 spec-grounded repair 연구에 특히 유용함을 시사한다.

\begin{table}[t]
\caption{Bug category distribution (3{,}691 AI bugs)}
\label{tab:taxonomy}
\centering
\begin{tabular}{lrr}
\toprule
Category & Count & Share \\
\midrule
Spec-Misinterpretation / Incomplete (WA) & 2{,}808 & 76\% \\
Inefficiency / Complexity (TLE)          & 504 & 14\% \\
Runtime error (index / input-format / numeric) & 338 & 9\% \\
Truncated / malformed output (CE)        & 41 & 1\% \\
\bottomrule
\end{tabular}
\end{table}

\section{Supported Tasks and Evaluation Protocols}
AIBugBench는 다음 다섯 태스크를 지원하며 각 태스크의 입력, 정답, 측정 지표를 표준화한다. Bug classification은 AI 버그를 결함 범주로 분류하는 태스크로, 라벨에 대한 accuracy와 두 평가자 간 Cohen's $\kappa$로 측정한다. Fault localization은 사람 버그의 fault location을 정답으로 사용하며 top-$k$ accuracy와 EXAM score로 측정한다~\cite{wong2016flsurvey}. Automatic program repair는 버그 프로그램과 전체 테스트, 그리고 참조 정답을 제공하며 수리 후 전체 테스트 통과율(plausible와 correct를 구분)로 측정한다~\cite{monperrus2018survey,xia2023llmapr}. AI 대 사람 버그 비교는 같은 문제에서 두 분포를 비교하는 태스크로, 범주 분포에 대한 $\chi^2$ 검정과 난이도별 교차표로 측정한다. Spec-grounded repair는 자연어 명세에서 추출한 제약을 두 번째 oracle로 활용하는 수리 평가로, 일반 수리 대비 Spec-Misinterpretation 한정 수리율 향상으로 측정한다.

표준 분할로는 문제 단위로 train/validation/test를 나누어 같은 문제의 솔루션이 분할을 넘나들지 않게 하여 누수 없는 평가를 보장한다. 모든 채점은 제공된 평가 스크립트와 동일한 시간 제한 및 라이브러리 환경에서 수행한다.

\section{Baseline Results}
생성 모델 baseline으로 \texttt{gpt-3.5-turbo}의 난이도별 통과율(표~\ref{tab:difficulty})을 보고한다. 전체 통과율은 13.8\%이며 어려운 문제에서 0\%로, 벤치마크가 충분히 도전적임을 보인다. 결함 분포(표~\ref{tab:taxonomy})는 결함 인식과 분류 태스크의 클래스 불균형 정도를 보여준다. 수리 baseline으로 교육용 수리 도구 Refactory~\cite{hu2019refactory}를 통합하여 범주별 수리율을 측정하는 프로토콜을 제공하며, 경쟁 프로그래밍 난이도와 도구의 입문용 가정 사이의 불일치가 기존 도구의 한계를 드러낼 것으로 예상한다. 전면적 수리 baseline 수치는 후속 확장에서 보고한다.

\section{Data Format, Access, and Maintenance}
데이터는 문제 단위 디렉터리로 제공된다. 각 문제 디렉터리는 정규화 명세(\texttt{problem.txt}), 전체 테스트(\texttt{testcases/}), AI 생성물(\texttt{ai/gen\_1..10.py}와 채점 라벨), 사람 프로그램(\texttt{human/}의 \texttt{faulty.py}, \texttt{correct.py}, \texttt{fault\-Location.txt}), 메타데이터(\texttt{meta.json})를 포함한다. 명세와 테스트, AI 생성물은 작은 텍스트로 배포되고, 대용량 hidden test는 ConDefects의 라이선스 하에서 재현 스크립트로 구축할 수 있다. 모든 코드와 평가 스크립트, manifest는 공개 저장소에서 버전 관리되며 새로운 모델로의 확장을 위한 생성과 채점 파이프라인을 함께 제공한다.

\section{Limitations and Ethics}
현재 AI 생성물은 단일 모델 \texttt{gpt-3.5-turbo}에 한정되며, 다양한 모델로의 확장은 제공된 파이프라인으로 가능하다. 사람 진정성은 ``컨테스트 시기가 ChatGPT 이전''이라는 proxy에 의존하며, 개별 제출 시각이 없으므로 옛 문제에 대한 후대 AI 재제출 가능성은 소수로 남는다. 일부 special judge 문제는 정확 일치 비교로 오판될 수 있어 메타데이터로 표시한다. 데이터는 AtCoder의 공개 문제와 ConDefects를 기반으로 하며, 비상업적 연구 목적의 공정 이용 범위에서 배포하고 원 출처와 라이선스를 명시한다.

\section{Conclusion}
AIBugBench는 모델 cutoff와 ChatGPT 출시 사이의 겹치는 구간을 이용하여, data leakage가 없고 사람과 같은 문제에서 짝지어진 LLM 생성 코드 결함 벤치마크를 제공한다. 정규화 명세, 전체 hidden test, fault location, AI와 사람 버그를 함께 담아 bug classification, fault localization, automatic program repair, AI 대 사람 비교, spec-grounded repair의 다섯 태스크를 지원한다. reference baseline으로 생성 모델의 난이도별 통과율과 결함 분포, systematicity를 보고하였으며, 향후 다중 모델 확장과 전면적 수리 baseline을 추가할 계획이다.

\section*{Data Availability}
데이터셋, 코드, 평가 스크립트는 \url{https://github.com/dwchoi95/aitaxo} 에서 공개한다.

\bibliographystyle{IEEEtran}
\bibliography{references}

\end{document}
